% ####################################################################
% Basic configuration of Beamera class and Jagiellonian theme
% ####################################################################
\RequirePackage[l2tabu, orthodox]{nag}



\ifx\PresentationStyle\notset
  \def\PresentationStyle{dark}
\fi



% Options: t - align frame text to the top.
\documentclass[10pt,t]{beamer}
\mode<presentation>
\usetheme[style=\PresentationStyle]{jagiellonian}




% ------------------------------------------------------------------------------------
% Procesing configuration files of Jagiellonian theme located
% in the directory "preambule"
% ------------------------------------------------------------------------------------
% Configuration for polish language
% Need description
\usepackage[polish]{babel}
% Need description
\usepackage[MeX]{polski}



% ------------------------------
% Better support of polish chars in technical parts of PDF
% ------------------------------
\hypersetup{pdfencoding=auto,psdextra}

% Package "textpos" give as enviroment "textblock" which is very usefull in
% arranging text on slides.

% This is standard configuration of "textpos"
\usepackage[overlay,absolute]{textpos}

% If you need to see bounds of "textblock's" comment line above and uncomment
% one below.

% Caution! When showboxes option is on significant ammunt of space is add
% to the top of textblock and as such, everyting put in them gone down.
% We need to check how to remove this bug.

% \usepackage[showboxes,overlay,absolute]{textpos}



% Setting scale length for package "textpos"
\setlength{\TPHorizModule}{10mm}
\setlength{\TPVertModule}{\TPHorizModule}


% ---------------------------------------
% Packages written for lectures "Geometria 3D dla twórców gier wideo"
% ---------------------------------------
% \usepackage{./Geometry3DPackages/Geometry3D}
% \usepackage{./Geometry3DPackages/Geometry3DIndices}
% \usepackage{./Geometry3DPackages/Geometry3DTikZStyle}
% \usepackage{./ProgramowanieSymulacjiFizykiPaczki/ProgramowanieSymulacjiFizykiTikZStyle}
% \usepackage{./Geometry3DPackages/mathcommands}


% ---------------------------------------
% Importing TikZ's libraries
% ---------------------------------------
\usetikzlibrary{arrows.meta}
\usetikzlibrary{positioning}





% % Configuration package "bm" that need for making bold symbols
\newcommand{\bmmax}{0}
\newcommand{\hmmax}{0}
\usepackage{bm}




% ---------------------------------------
% Packages for scientific texts
% ---------------------------------------
% \let\lll\undefined  % Sometimes you must use this line to allow
% "amsmath" package to works with packages with packages for polish
% languge imported
% /preambul/LanguageSettings/JagiellonianPolishLanguageSettings.tex.
% This comments (probably) removes polish letter Ł.
\usepackage{amsmath}  % Packages from American Mathematical Society (AMS)
\usepackage{amssymb}
\usepackage{amscd}
\usepackage{amsthm}
\usepackage{siunitx}  % Package for typsetting SI units.
\usepackage{upgreek}  % Better looking greek letters.
% Example of using upgreek: pi = \uppi


\usepackage{calrsfs}  % Zmienia czcionkę kaligraficzną w \mathcal
% na ładniejszą. Może w innych miejscach robi to samo, ale o tym nic
% nie wiem.







% !!!!!!!!!!!!!!!!!!!!!!!!!!!!!!
% !!!!!!!!!!!!!!!!!!!!!!!!!!!!!!
% EVIL STUFF
\if\JUlogotitle1
\edef\LogoJUPath{LogoJU_\JUlogoLang/LogoJU_\JUlogoShape_\JUlogoColor.pdf}
\titlegraphic{\hfill\includegraphics[scale=0.22]
{./JagiellonianPictures/\LogoJUPath}}
\fi

% ---------------------------------------
% Commands for handling colors
% ---------------------------------------


% Command for setting normal text color for some text in math modestyle
% Text color depend on used style of Jagiellonian

% Beamer version of command
\newcommand{\TextWithNormalTextColor}[1]{%
  {\color{jNormalTextFGColor}
    \setbeamercolor{math text}{fg=jNormalTextFGColor} {#1}}
}

% Article and similar classes version of command
% \newcommand{\TextWithNormalTextColor}[1]{%
%   {\color{jNormalTextsFGColor} {#1}}
% }



% Beamer version of command
\newcommand{\NormalTextInMathMode}[1]{%
  {\color{jNormalTextFGColor}
    \setbeamercolor{math text}{fg=jNormalTextFGColor} \text{#1}}
}


% Article and similar classes version of command
% \newcommand{\NormalTextInMathMode}[1]{%
%   {\color{jNormalTextsFGColor} \text{#1}}
% }




% Command that sets color of some mathematical text to the same color
% that has normal text in header (?)

% Beamer version of the command
\newcommand{\MathTextFrametitleFGColor}[1]{%
  {\color{jFrametitleFGColor}
    \setbeamercolor{math text}{fg=jFrametitleFGColor} #1}
}

% Article and similar classes version of the command
% \newcommand{\MathTextWhiteColor}[1]{{\color{jFrametitleFGColor} #1}}





% Command for setting color of alert text for some text in math modestyle

% Beamer version of the command
\newcommand{\MathTextAlertColor}[1]{%
  {\color{jOrange} \setbeamercolor{math text}{fg=jOrange} #1}
}

% Article and similar classes version of the command
% \newcommand{\MathTextAlertColor}[1]{{\color{jOrange} #1}}





% Command that allow you to sets chosen color as the color of some text into
% math mode. Due to some nuances in the way that Beamer handle colors
% it not work in all cases. We hope that in the future we will improve it.

% Beamer version of the command
\newcommand{\SetMathTextColor}[2]{%
  {\color{#1} \setbeamercolor{math text}{fg=#1} #2}
}


% Article and similar classes version of the command
% \newcommand{\SetMathTextColor}[2]{{\color{#1} #2}}










% ---------------------------------------
% Commands for few special slides
% ---------------------------------------
\newcommand{\EndingSlide}[1]{%
  \begin{frame}[standout]

    \begingroup

    \color{jFrametitleFGColor}

    #1

    \endgroup

  \end{frame}
}










% ---------------------------------------
% Commands for setting background pictures for some slides
% ---------------------------------------
\newcommand{\TitleBackgroundPicture}
{./JagiellonianPictures/Backgrounds/LajkonikDark.png}
\newcommand{\SectionBackgroundPicture}
{./JagiellonianPictures/Backgrounds/LajkonikLight.png}



\newcommand{\TitleSlideWithPicture}{%
  \begingroup

  \usebackgroundtemplate{%
    \includegraphics[height=\paperheight]{\TitleBackgroundPicture}}

  \maketitle

  \endgroup
}





\newcommand{\SectionSlideWithPicture}[1]{%
  \begingroup

  \usebackgroundtemplate{%
    \includegraphics[height=\paperheight]{\SectionBackgroundPicture}}

  \setbeamercolor{titlelike}{fg=normal text.fg}

  \section{#1}

  \endgroup
}










% ---------------------------------------
% Commands for lectures "Geometria 3D dla twórców gier wideo"
% Polish version
% ---------------------------------------
% Komendy teraz wykomentowane były potrzebne, gdy loga były na niebieskim
% tle, nie na białym. A są na białym bo tego chcieli w biurze projektu.
% \newcommand{\FundingLogoWhitePicturePL}
% {./PresentationPictures/CommonPictures/logotypFundusze_biale_bez_tla2.pdf}
\newcommand{\FundingLogoColorPicturePL}
{./PresentationPictures/CommonPictures/European_Funds_color_PL.pdf}
% \newcommand{\EULogoWhitePicturePL}
% {./PresentationPictures/CommonPictures/logotypUE_biale_bez_tla2.pdf}
\newcommand{\EUSocialFundLogoColorPicturePL}
{./PresentationPictures/CommonPictures/EU_Social_Fund_color_PL.pdf}
% \newcommand{\ZintegrUJLogoWhitePicturePL}
% {./PresentationPictures/CommonPictures/zintegruj-logo-white.pdf}
\newcommand{\ZintegrUJLogoColorPicturePL}
{./PresentationPictures/CommonPictures/ZintegrUJ_color.pdf}
\newcommand{\JULogoColorPicturePL}
{./JagiellonianPictures/LogoJU_PL/LogoJU_A_color.pdf}





\newcommand{\GeometryThreeDSpecialBeginningSlidePL}{%
  \begin{frame}[standout]

    \begin{textblock}{11}(1,0.7)

      \begin{flushleft}

        \mdseries

        \footnotesize

        \color{jFrametitleFGColor}

        Materiał powstał w ramach projektu współfinansowanego ze środków
        Unii Europejskiej w ramach Europejskiego Funduszu Społecznego
        POWR.03.05.00-00-Z309/17-00.

      \end{flushleft}

    \end{textblock}





    \begin{textblock}{10}(0,2.2)

      \tikz \fill[color=jBackgroundStyleLight] (0,0) rectangle (12.8,-1.5);

    \end{textblock}


    \begin{textblock}{3.2}(1,2.45)

      \includegraphics[scale=0.3]{\FundingLogoColorPicturePL}

    \end{textblock}


    \begin{textblock}{2.5}(3.7,2.5)

      \includegraphics[scale=0.2]{\JULogoColorPicturePL}

    \end{textblock}


    \begin{textblock}{2.5}(6,2.4)

      \includegraphics[scale=0.1]{\ZintegrUJLogoColorPicturePL}

    \end{textblock}


    \begin{textblock}{4.2}(8.4,2.6)

      \includegraphics[scale=0.3]{\EUSocialFundLogoColorPicturePL}

    \end{textblock}

  \end{frame}
}



\newcommand{\GeometryThreeDTwoSpecialBeginningSlidesPL}{%
  \begin{frame}[standout]

    \begin{textblock}{11}(1,0.7)

      \begin{flushleft}

        \mdseries

        \footnotesize

        \color{jFrametitleFGColor}

        Materiał powstał w ramach projektu współfinansowanego ze środków
        Unii Europejskiej w ramach Europejskiego Funduszu Społecznego
        POWR.03.05.00-00-Z309/17-00.

      \end{flushleft}

    \end{textblock}





    \begin{textblock}{10}(0,2.2)

      \tikz \fill[color=jBackgroundStyleLight] (0,0) rectangle (12.8,-1.5);

    \end{textblock}


    \begin{textblock}{3.2}(1,2.45)

      \includegraphics[scale=0.3]{\FundingLogoColorPicturePL}

    \end{textblock}


    \begin{textblock}{2.5}(3.7,2.5)

      \includegraphics[scale=0.2]{\JULogoColorPicturePL}

    \end{textblock}


    \begin{textblock}{2.5}(6,2.4)

      \includegraphics[scale=0.1]{\ZintegrUJLogoColorPicturePL}

    \end{textblock}


    \begin{textblock}{4.2}(8.4,2.6)

      \includegraphics[scale=0.3]{\EUSocialFundLogoColorPicturePL}

    \end{textblock}

  \end{frame}





  \TitleSlideWithPicture
}



\newcommand{\GeometryThreeDSpecialEndingSlidePL}{%
  \begin{frame}[standout]

    \begin{textblock}{11}(1,0.7)

      \begin{flushleft}

        \mdseries

        \footnotesize

        \color{jFrametitleFGColor}

        Materiał powstał w ramach projektu współfinansowanego ze środków
        Unii Europejskiej w~ramach Europejskiego Funduszu Społecznego
        POWR.03.05.00-00-Z309/17-00.

      \end{flushleft}

    \end{textblock}





    \begin{textblock}{10}(0,2.2)

      \tikz \fill[color=jBackgroundStyleLight] (0,0) rectangle (12.8,-1.5);

    \end{textblock}


    \begin{textblock}{3.2}(1,2.45)

      \includegraphics[scale=0.3]{\FundingLogoColorPicturePL}

    \end{textblock}


    \begin{textblock}{2.5}(3.7,2.5)

      \includegraphics[scale=0.2]{\JULogoColorPicturePL}

    \end{textblock}


    \begin{textblock}{2.5}(6,2.4)

      \includegraphics[scale=0.1]{\ZintegrUJLogoColorPicturePL}

    \end{textblock}


    \begin{textblock}{4.2}(8.4,2.6)

      \includegraphics[scale=0.3]{\EUSocialFundLogoColorPicturePL}

    \end{textblock}





    \begin{textblock}{11}(1,4)

      \begin{flushleft}

        \mdseries

        \footnotesize

        \RaggedRight

        \color{jFrametitleFGColor}

        Treść niniejszego wykładu jest udostępniona na~licencji
        Creative Commons (\textsc{cc}), z~uzna\-niem autorstwa
        (\textsc{by}) oraz udostępnianiem na tych samych warunkach
        (\textsc{sa}). Rysunki i~wy\-kresy zawarte w~wykładzie są
        autorstwa dr.~hab.~Pawła Węgrzyna et~al. i~są dostępne
        na tej samej licencji, o~ile nie wskazano inaczej.
        W~prezentacji wykorzystano temat Beamera Jagiellonian,
        oparty na~temacie Metropolis Matthiasa Vogelgesanga,
        dostępnym na licencji \LaTeX{} Project Public License~1.3c
        pod adresem: \colorhref{https://github.com/matze/mtheme}
        {https://github.com/matze/mtheme}.

        Projekt typograficzny: Iwona Grabska-Gradzińska \\
        Skład: Kamil Ziemian;
        Korekta: Wojciech Palacz \\
        Modele: Dariusz Frymus, Kamil Nowakowski \\
        Rysunki i~wykresy: Kamil Ziemian, Paweł Węgrzyn, Wojciech Palacz

      \end{flushleft}

    \end{textblock}

  \end{frame}
}



\newcommand{\GeometryThreeDTwoSpecialEndingSlidesPL}[1]{%
  \begin{frame}[standout]


    \begin{textblock}{11}(1,0.7)

      \begin{flushleft}

        \mdseries

        \footnotesize

        \color{jFrametitleFGColor}

        Materiał powstał w ramach projektu współfinansowanego ze środków
        Unii Europejskiej w~ramach Europejskiego Funduszu Społecznego
        POWR.03.05.00-00-Z309/17-00.

      \end{flushleft}

    \end{textblock}





    \begin{textblock}{10}(0,2.2)

      \tikz \fill[color=jBackgroundStyleLight] (0,0) rectangle (12.8,-1.5);

    \end{textblock}


    \begin{textblock}{3.2}(1,2.45)

      \includegraphics[scale=0.3]{\FundingLogoColorPicturePL}

    \end{textblock}


    \begin{textblock}{2.5}(3.7,2.5)

      \includegraphics[scale=0.2]{\JULogoColorPicturePL}

    \end{textblock}


    \begin{textblock}{2.5}(6,2.4)

      \includegraphics[scale=0.1]{\ZintegrUJLogoColorPicturePL}

    \end{textblock}


    \begin{textblock}{4.2}(8.4,2.6)

      \includegraphics[scale=0.3]{\EUSocialFundLogoColorPicturePL}

    \end{textblock}





    \begin{textblock}{11}(1,4)

      \begin{flushleft}

        \mdseries

        \footnotesize

        \RaggedRight

        \color{jFrametitleFGColor}

        Treść niniejszego wykładu jest udostępniona na~licencji
        Creative Commons (\textsc{cc}), z~uzna\-niem autorstwa
        (\textsc{by}) oraz udostępnianiem na tych samych warunkach
        (\textsc{sa}). Rysunki i~wy\-kresy zawarte w~wykładzie są
        autorstwa dr.~hab.~Pawła Węgrzyna et~al. i~są dostępne
        na tej samej licencji, o~ile nie wskazano inaczej.
        W~prezentacji wykorzystano temat Beamera Jagiellonian,
        oparty na~temacie Metropolis Matthiasa Vogelgesanga,
        dostępnym na licencji \LaTeX{} Project Public License~1.3c
        pod adresem: \colorhref{https://github.com/matze/mtheme}
        {https://github.com/matze/mtheme}.

        Projekt typograficzny: Iwona Grabska-Gradzińska \\
        Skład: Kamil Ziemian;
        Korekta: Wojciech Palacz \\
        Modele: Dariusz Frymus, Kamil Nowakowski \\
        Rysunki i~wykresy: Kamil Ziemian, Paweł Węgrzyn, Wojciech Palacz

      \end{flushleft}

    \end{textblock}

  \end{frame}





  \begin{frame}[standout]

    \begingroup

    \color{jFrametitleFGColor}

    #1

    \endgroup

  \end{frame}
}



\newcommand{\GeometryThreeDSpecialEndingSlideVideoPL}{%
  \begin{frame}[standout]

    \begin{textblock}{11}(1,0.7)

      \begin{flushleft}

        \mdseries

        \footnotesize

        \color{jFrametitleFGColor}

        Materiał powstał w ramach projektu współfinansowanego ze środków
        Unii Europejskiej w~ramach Europejskiego Funduszu Społecznego
        POWR.03.05.00-00-Z309/17-00.

      \end{flushleft}

    \end{textblock}





    \begin{textblock}{10}(0,2.2)

      \tikz \fill[color=jBackgroundStyleLight] (0,0) rectangle (12.8,-1.5);

    \end{textblock}


    \begin{textblock}{3.2}(1,2.45)

      \includegraphics[scale=0.3]{\FundingLogoColorPicturePL}

    \end{textblock}


    \begin{textblock}{2.5}(3.7,2.5)

      \includegraphics[scale=0.2]{\JULogoColorPicturePL}

    \end{textblock}


    \begin{textblock}{2.5}(6,2.4)

      \includegraphics[scale=0.1]{\ZintegrUJLogoColorPicturePL}

    \end{textblock}


    \begin{textblock}{4.2}(8.4,2.6)

      \includegraphics[scale=0.3]{\EUSocialFundLogoColorPicturePL}

    \end{textblock}





    \begin{textblock}{11}(1,4)

      \begin{flushleft}

        \mdseries

        \footnotesize

        \RaggedRight

        \color{jFrametitleFGColor}

        Treść niniejszego wykładu jest udostępniona na~licencji
        Creative Commons (\textsc{cc}), z~uzna\-niem autorstwa
        (\textsc{by}) oraz udostępnianiem na tych samych warunkach
        (\textsc{sa}). Rysunki i~wy\-kresy zawarte w~wykładzie są
        autorstwa dr.~hab.~Pawła Węgrzyna et~al. i~są dostępne
        na tej samej licencji, o~ile nie wskazano inaczej.
        W~prezentacji wykorzystano temat Beamera Jagiellonian,
        oparty na~temacie Metropolis Matthiasa Vogelgesanga,
        dostępnym na licencji \LaTeX{} Project Public License~1.3c
        pod adresem: \colorhref{https://github.com/matze/mtheme}
        {https://github.com/matze/mtheme}.

        Projekt typograficzny: Iwona Grabska-Gradzińska;
        Skład: Kamil Ziemian \\
        Korekta: Wojciech Palacz;
        Modele: Dariusz Frymus, Kamil Nowakowski \\
        Rysunki i~wykresy: Kamil Ziemian, Paweł Węgrzyn, Wojciech Palacz \\
        Montaż: Agencja Filmowa Film \& Television Production~-- Zbigniew
        Masklak

      \end{flushleft}

    \end{textblock}

  \end{frame}
}





\newcommand{\GeometryThreeDTwoSpecialEndingSlidesVideoPL}[1]{%
  \begin{frame}[standout]

    \begin{textblock}{11}(1,0.7)

      \begin{flushleft}

        \mdseries

        \footnotesize

        \color{jFrametitleFGColor}

        Materiał powstał w ramach projektu współfinansowanego ze środków
        Unii Europejskiej w~ramach Europejskiego Funduszu Społecznego
        POWR.03.05.00-00-Z309/17-00.

      \end{flushleft}

    \end{textblock}





    \begin{textblock}{10}(0,2.2)

      \tikz \fill[color=jBackgroundStyleLight] (0,0) rectangle (12.8,-1.5);

    \end{textblock}


    \begin{textblock}{3.2}(1,2.45)

      \includegraphics[scale=0.3]{\FundingLogoColorPicturePL}

    \end{textblock}


    \begin{textblock}{2.5}(3.7,2.5)

      \includegraphics[scale=0.2]{\JULogoColorPicturePL}

    \end{textblock}


    \begin{textblock}{2.5}(6,2.4)

      \includegraphics[scale=0.1]{\ZintegrUJLogoColorPicturePL}

    \end{textblock}


    \begin{textblock}{4.2}(8.4,2.6)

      \includegraphics[scale=0.3]{\EUSocialFundLogoColorPicturePL}

    \end{textblock}





    \begin{textblock}{11}(1,4)

      \begin{flushleft}

        \mdseries

        \footnotesize

        \RaggedRight

        \color{jFrametitleFGColor}

        Treść niniejszego wykładu jest udostępniona na~licencji
        Creative Commons (\textsc{cc}), z~uzna\-niem autorstwa
        (\textsc{by}) oraz udostępnianiem na tych samych warunkach
        (\textsc{sa}). Rysunki i~wy\-kresy zawarte w~wykładzie są
        autorstwa dr.~hab.~Pawła Węgrzyna et~al. i~są dostępne
        na tej samej licencji, o~ile nie wskazano inaczej.
        W~prezentacji wykorzystano temat Beamera Jagiellonian,
        oparty na~temacie Metropolis Matthiasa Vogelgesanga,
        dostępnym na licencji \LaTeX{} Project Public License~1.3c
        pod adresem: \colorhref{https://github.com/matze/mtheme}
        {https://github.com/matze/mtheme}.

        Projekt typograficzny: Iwona Grabska-Gradzińska;
        Skład: Kamil Ziemian \\
        Korekta: Wojciech Palacz;
        Modele: Dariusz Frymus, Kamil Nowakowski \\
        Rysunki i~wykresy: Kamil Ziemian, Paweł Węgrzyn, Wojciech Palacz \\
        Montaż: Agencja Filmowa Film \& Television Production~-- Zbigniew
        Masklak

      \end{flushleft}

    \end{textblock}

  \end{frame}





  \begin{frame}[standout]


    \begingroup

    \color{jFrametitleFGColor}

    #1

    \endgroup

  \end{frame}
}










% ---------------------------------------
% Commands for lectures "Geometria 3D dla twórców gier wideo"
% English version
% ---------------------------------------
% \newcommand{\FundingLogoWhitePictureEN}
% {./PresentationPictures/CommonPictures/logotypFundusze_biale_bez_tla2.pdf}
\newcommand{\FundingLogoColorPictureEN}
{./PresentationPictures/CommonPictures/European_Funds_color_EN.pdf}
% \newcommand{\EULogoWhitePictureEN}
% {./PresentationPictures/CommonPictures/logotypUE_biale_bez_tla2.pdf}
\newcommand{\EUSocialFundLogoColorPictureEN}
{./PresentationPictures/CommonPictures/EU_Social_Fund_color_EN.pdf}
% \newcommand{\ZintegrUJLogoWhitePictureEN}
% {./PresentationPictures/CommonPictures/zintegruj-logo-white.pdf}
\newcommand{\ZintegrUJLogoColorPictureEN}
{./PresentationPictures/CommonPictures/ZintegrUJ_color.pdf}
\newcommand{\JULogoColorPictureEN}
{./JagiellonianPictures/LogoJU_EN/LogoJU_A_color.pdf}



\newcommand{\GeometryThreeDSpecialBeginningSlideEN}{%
  \begin{frame}[standout]

    \begin{textblock}{11}(1,0.7)

      \begin{flushleft}

        \mdseries

        \footnotesize

        \color{jFrametitleFGColor}

        This content was created as part of a project co-financed by the
        European Union within the framework of the European Social Fund
        POWR.03.05.00-00-Z309/17-00.

      \end{flushleft}

    \end{textblock}





    \begin{textblock}{10}(0,2.2)

      \tikz \fill[color=jBackgroundStyleLight] (0,0) rectangle (12.8,-1.5);

    \end{textblock}


    \begin{textblock}{3.2}(0.7,2.45)

      \includegraphics[scale=0.3]{\FundingLogoColorPictureEN}

    \end{textblock}


    \begin{textblock}{2.5}(4.15,2.5)

      \includegraphics[scale=0.2]{\JULogoColorPictureEN}

    \end{textblock}


    \begin{textblock}{2.5}(6.35,2.4)

      \includegraphics[scale=0.1]{\ZintegrUJLogoColorPictureEN}

    \end{textblock}


    \begin{textblock}{4.2}(8.4,2.6)

      \includegraphics[scale=0.3]{\EUSocialFundLogoColorPictureEN}

    \end{textblock}

  \end{frame}
}



\newcommand{\GeometryThreeDTwoSpecialBeginningSlidesEN}{%
  \begin{frame}[standout]

    \begin{textblock}{11}(1,0.7)

      \begin{flushleft}

        \mdseries

        \footnotesize

        \color{jFrametitleFGColor}

        This content was created as part of a project co-financed by the
        European Union within the framework of the European Social Fund
        POWR.03.05.00-00-Z309/17-00.

      \end{flushleft}

    \end{textblock}





    \begin{textblock}{10}(0,2.2)

      \tikz \fill[color=jBackgroundStyleLight] (0,0) rectangle (12.8,-1.5);

    \end{textblock}


    \begin{textblock}{3.2}(0.7,2.45)

      \includegraphics[scale=0.3]{\FundingLogoColorPictureEN}

    \end{textblock}


    \begin{textblock}{2.5}(4.15,2.5)

      \includegraphics[scale=0.2]{\JULogoColorPictureEN}

    \end{textblock}


    \begin{textblock}{2.5}(6.35,2.4)

      \includegraphics[scale=0.1]{\ZintegrUJLogoColorPictureEN}

    \end{textblock}


    \begin{textblock}{4.2}(8.4,2.6)

      \includegraphics[scale=0.3]{\EUSocialFundLogoColorPictureEN}

    \end{textblock}

  \end{frame}





  \TitleSlideWithPicture
}



\newcommand{\GeometryThreeDSpecialEndingSlideEN}{%
  \begin{frame}[standout]

    \begin{textblock}{11}(1,0.7)

      \begin{flushleft}

        \mdseries

        \footnotesize

        \color{jFrametitleFGColor}

        This content was created as part of a project co-financed by the
        European Union within the framework of the European Social Fund
        POWR.03.05.00-00-Z309/17-00.

      \end{flushleft}

    \end{textblock}





    \begin{textblock}{10}(0,2.2)

      \tikz \fill[color=jBackgroundStyleLight] (0,0) rectangle (12.8,-1.5);

    \end{textblock}


    \begin{textblock}{3.2}(0.7,2.45)

      \includegraphics[scale=0.3]{\FundingLogoColorPictureEN}

    \end{textblock}


    \begin{textblock}{2.5}(4.15,2.5)

      \includegraphics[scale=0.2]{\JULogoColorPictureEN}

    \end{textblock}


    \begin{textblock}{2.5}(6.35,2.4)

      \includegraphics[scale=0.1]{\ZintegrUJLogoColorPictureEN}

    \end{textblock}


    \begin{textblock}{4.2}(8.4,2.6)

      \includegraphics[scale=0.3]{\EUSocialFundLogoColorPictureEN}

    \end{textblock}





    \begin{textblock}{11}(1,4)

      \begin{flushleft}

        \mdseries

        \footnotesize

        \RaggedRight

        \color{jFrametitleFGColor}

        The content of this lecture is made available under a~Creative
        Commons licence (\textsc{cc}), giving the author the credits
        (\textsc{by}) and putting an obligation to share on the same terms
        (\textsc{sa}). Figures and diagrams included in the lecture are
        authored by Paweł Węgrzyn et~al., and are available under the same
        license unless indicated otherwise.\\ The presentation uses the
        Beamer Jagiellonian theme based on Matthias Vogelgesang’s
        Metropolis theme, available under license \LaTeX{} Project
        Public License~1.3c at: \colorhref{https://github.com/matze/mtheme}
        {https://github.com/matze/mtheme}.

        Typographic design: Iwona Grabska-Gradzińska \\
        \LaTeX{} Typesetting: Kamil Ziemian \\
        Proofreading: Wojciech Palacz,
        Monika Stawicka \\
        3D Models: Dariusz Frymus, Kamil Nowakowski \\
        Figures and charts: Kamil Ziemian, Paweł Węgrzyn, Wojciech Palacz

      \end{flushleft}

    \end{textblock}

  \end{frame}
}



\newcommand{\GeometryThreeDTwoSpecialEndingSlidesEN}[1]{%
  \begin{frame}[standout]


    \begin{textblock}{11}(1,0.7)

      \begin{flushleft}

        \mdseries

        \footnotesize

        \color{jFrametitleFGColor}

        This content was created as part of a project co-financed by the
        European Union within the framework of the European Social Fund
        POWR.03.05.00-00-Z309/17-00.

      \end{flushleft}

    \end{textblock}





    \begin{textblock}{10}(0,2.2)

      \tikz \fill[color=jBackgroundStyleLight] (0,0) rectangle (12.8,-1.5);

    \end{textblock}


    \begin{textblock}{3.2}(0.7,2.45)

      \includegraphics[scale=0.3]{\FundingLogoColorPictureEN}

    \end{textblock}


    \begin{textblock}{2.5}(4.15,2.5)

      \includegraphics[scale=0.2]{\JULogoColorPictureEN}

    \end{textblock}


    \begin{textblock}{2.5}(6.35,2.4)

      \includegraphics[scale=0.1]{\ZintegrUJLogoColorPictureEN}

    \end{textblock}


    \begin{textblock}{4.2}(8.4,2.6)

      \includegraphics[scale=0.3]{\EUSocialFundLogoColorPictureEN}

    \end{textblock}





    \begin{textblock}{11}(1,4)

      \begin{flushleft}

        \mdseries

        \footnotesize

        \RaggedRight

        \color{jFrametitleFGColor}

        The content of this lecture is made available under a~Creative
        Commons licence (\textsc{cc}), giving the author the credits
        (\textsc{by}) and putting an obligation to share on the same terms
        (\textsc{sa}). Figures and diagrams included in the lecture are
        authored by Paweł Węgrzyn et~al., and are available under the same
        license unless indicated otherwise.\\ The presentation uses the
        Beamer Jagiellonian theme based on Matthias Vogelgesang’s
        Metropolis theme, available under license \LaTeX{} Project
        Public License~1.3c at: \colorhref{https://github.com/matze/mtheme}
        {https://github.com/matze/mtheme}.

        Typographic design: Iwona Grabska-Gradzińska \\
        \LaTeX{} Typesetting: Kamil Ziemian \\
        Proofreading: Wojciech Palacz,
        Monika Stawicka \\
        3D Models: Dariusz Frymus, Kamil Nowakowski \\
        Figures and charts: Kamil Ziemian, Paweł Węgrzyn, Wojciech Palacz

      \end{flushleft}

    \end{textblock}

  \end{frame}





  \begin{frame}[standout]

    \begingroup

    \color{jFrametitleFGColor}

    #1

    \endgroup

  \end{frame}
}



\newcommand{\GeometryThreeDSpecialEndingSlideVideoVerOneEN}{%
  \begin{frame}[standout]

    \begin{textblock}{11}(1,0.7)

      \begin{flushleft}

        \mdseries

        \footnotesize

        \color{jFrametitleFGColor}

        This content was created as part of a project co-financed by the
        European Union within the framework of the European Social Fund
        POWR.03.05.00-00-Z309/17-00.

      \end{flushleft}

    \end{textblock}





    \begin{textblock}{10}(0,2.2)

      \tikz \fill[color=jBackgroundStyleLight] (0,0) rectangle (12.8,-1.5);

    \end{textblock}


    \begin{textblock}{3.2}(0.7,2.45)

      \includegraphics[scale=0.3]{\FundingLogoColorPictureEN}

    \end{textblock}


    \begin{textblock}{2.5}(4.15,2.5)

      \includegraphics[scale=0.2]{\JULogoColorPictureEN}

    \end{textblock}


    \begin{textblock}{2.5}(6.35,2.4)

      \includegraphics[scale=0.1]{\ZintegrUJLogoColorPictureEN}

    \end{textblock}


    \begin{textblock}{4.2}(8.4,2.6)

      \includegraphics[scale=0.3]{\EUSocialFundLogoColorPictureEN}

    \end{textblock}





    \begin{textblock}{11}(1,4)

      \begin{flushleft}

        \mdseries

        \footnotesize

        \RaggedRight

        \color{jFrametitleFGColor}

        The content of this lecture is made available under a Creative
        Commons licence (\textsc{cc}), giving the author the credits
        (\textsc{by}) and putting an obligation to share on the same terms
        (\textsc{sa}). Figures and diagrams included in the lecture are
        authored by Paweł Węgrzyn et~al., and are available under the same
        license unless indicated otherwise.\\ The presentation uses the
        Beamer Jagiellonian theme based on Matthias Vogelgesang’s
        Metropolis theme, available under license \LaTeX{} Project
        Public License~1.3c at: \colorhref{https://github.com/matze/mtheme}
        {https://github.com/matze/mtheme}.

        Typographic design: Iwona Grabska-Gradzińska;
        \LaTeX{} Typesetting: Kamil Ziemian \\
        Proofreading: Wojciech Palacz,
        Monika Stawicka \\
        3D Models: Dariusz Frymus, Kamil Nowakowski \\
        Figures and charts: Kamil Ziemian, Paweł Węgrzyn, Wojciech
        Palacz \\
        Film editing: Agencja Filmowa Film \& Television Production~--
        Zbigniew Masklak

      \end{flushleft}

    \end{textblock}

  \end{frame}
}



\newcommand{\GeometryThreeDSpecialEndingSlideVideoVerTwoEN}{%
  \begin{frame}[standout]

    \begin{textblock}{11}(1,0.7)

      \begin{flushleft}

        \mdseries

        \footnotesize

        \color{jFrametitleFGColor}

        This content was created as part of a project co-financed by the
        European Union within the framework of the European Social Fund
        POWR.03.05.00-00-Z309/17-00.

      \end{flushleft}

    \end{textblock}





    \begin{textblock}{10}(0,2.2)

      \tikz \fill[color=jBackgroundStyleLight] (0,0) rectangle (12.8,-1.5);

    \end{textblock}


    \begin{textblock}{3.2}(0.7,2.45)

      \includegraphics[scale=0.3]{\FundingLogoColorPictureEN}

    \end{textblock}


    \begin{textblock}{2.5}(4.15,2.5)

      \includegraphics[scale=0.2]{\JULogoColorPictureEN}

    \end{textblock}


    \begin{textblock}{2.5}(6.35,2.4)

      \includegraphics[scale=0.1]{\ZintegrUJLogoColorPictureEN}

    \end{textblock}


    \begin{textblock}{4.2}(8.4,2.6)

      \includegraphics[scale=0.3]{\EUSocialFundLogoColorPictureEN}

    \end{textblock}





    \begin{textblock}{11}(1,4)

      \begin{flushleft}

        \mdseries

        \footnotesize

        \RaggedRight

        \color{jFrametitleFGColor}

        The content of this lecture is made available under a Creative
        Commons licence (\textsc{cc}), giving the author the credits
        (\textsc{by}) and putting an obligation to share on the same terms
        (\textsc{sa}). Figures and diagrams included in the lecture are
        authored by Paweł Węgrzyn et~al., and are available under the same
        license unless indicated otherwise.\\ The presentation uses the
        Beamer Jagiellonian theme based on Matthias Vogelgesang’s
        Metropolis theme, available under license \LaTeX{} Project
        Public License~1.3c at: \colorhref{https://github.com/matze/mtheme}
        {https://github.com/matze/mtheme}.

        Typographic design: Iwona Grabska-Gradzińska;
        \LaTeX{} Typesetting: Kamil Ziemian \\
        Proofreading: Wojciech Palacz,
        Monika Stawicka \\
        3D Models: Dariusz Frymus, Kamil Nowakowski \\
        Figures and charts: Kamil Ziemian, Paweł Węgrzyn, Wojciech
        Palacz \\
        Film editing: IMAVI -- Joanna Kozakiewicz, Krzysztof Magda, Nikodem
        Frodyma

      \end{flushleft}

    \end{textblock}

  \end{frame}
}



\newcommand{\GeometryThreeDSpecialEndingSlideVideoVerThreeEN}{%
  \begin{frame}[standout]

    \begin{textblock}{11}(1,0.7)

      \begin{flushleft}

        \mdseries

        \footnotesize

        \color{jFrametitleFGColor}

        This content was created as part of a project co-financed by the
        European Union within the framework of the European Social Fund
        POWR.03.05.00-00-Z309/17-00.

      \end{flushleft}

    \end{textblock}





    \begin{textblock}{10}(0,2.2)

      \tikz \fill[color=jBackgroundStyleLight] (0,0) rectangle (12.8,-1.5);

    \end{textblock}


    \begin{textblock}{3.2}(0.7,2.45)

      \includegraphics[scale=0.3]{\FundingLogoColorPictureEN}

    \end{textblock}


    \begin{textblock}{2.5}(4.15,2.5)

      \includegraphics[scale=0.2]{\JULogoColorPictureEN}

    \end{textblock}


    \begin{textblock}{2.5}(6.35,2.4)

      \includegraphics[scale=0.1]{\ZintegrUJLogoColorPictureEN}

    \end{textblock}


    \begin{textblock}{4.2}(8.4,2.6)

      \includegraphics[scale=0.3]{\EUSocialFundLogoColorPictureEN}

    \end{textblock}





    \begin{textblock}{11}(1,4)

      \begin{flushleft}

        \mdseries

        \footnotesize

        \RaggedRight

        \color{jFrametitleFGColor}

        The content of this lecture is made available under a Creative
        Commons licence (\textsc{cc}), giving the author the credits
        (\textsc{by}) and putting an obligation to share on the same terms
        (\textsc{sa}). Figures and diagrams included in the lecture are
        authored by Paweł Węgrzyn et~al., and are available under the same
        license unless indicated otherwise.\\ The presentation uses the
        Beamer Jagiellonian theme based on Matthias Vogelgesang’s
        Metropolis theme, available under license \LaTeX{} Project
        Public License~1.3c at: \colorhref{https://github.com/matze/mtheme}
        {https://github.com/matze/mtheme}.

        Typographic design: Iwona Grabska-Gradzińska;
        \LaTeX{} Typesetting: Kamil Ziemian \\
        Proofreading: Wojciech Palacz,
        Monika Stawicka \\
        3D Models: Dariusz Frymus, Kamil Nowakowski \\
        Figures and charts: Kamil Ziemian, Paweł Węgrzyn, Wojciech
        Palacz \\
        Film editing: Agencja Filmowa Film \& Television Production~--
        Zbigniew Masklak \\
        Film editing: IMAVI -- Joanna Kozakiewicz, Krzysztof Magda, Nikodem
        Frodyma

      \end{flushleft}

    \end{textblock}

  \end{frame}
}



\newcommand{\GeometryThreeDTwoSpecialEndingSlidesVideoVerOneEN}[1]{%
  \begin{frame}[standout]

    \begin{textblock}{11}(1,0.7)

      \begin{flushleft}

        \mdseries

        \footnotesize

        \color{jFrametitleFGColor}

        This content was created as part of a project co-financed by the
        European Union within the framework of the European Social Fund
        POWR.03.05.00-00-Z309/17-00.

      \end{flushleft}

    \end{textblock}





    \begin{textblock}{10}(0,2.2)

      \tikz \fill[color=jBackgroundStyleLight] (0,0) rectangle (12.8,-1.5);

    \end{textblock}


    \begin{textblock}{3.2}(0.7,2.45)

      \includegraphics[scale=0.3]{\FundingLogoColorPictureEN}

    \end{textblock}


    \begin{textblock}{2.5}(4.15,2.5)

      \includegraphics[scale=0.2]{\JULogoColorPictureEN}

    \end{textblock}


    \begin{textblock}{2.5}(6.35,2.4)

      \includegraphics[scale=0.1]{\ZintegrUJLogoColorPictureEN}

    \end{textblock}


    \begin{textblock}{4.2}(8.4,2.6)

      \includegraphics[scale=0.3]{\EUSocialFundLogoColorPictureEN}

    \end{textblock}





    \begin{textblock}{11}(1,4)

      \begin{flushleft}

        \mdseries

        \footnotesize

        \RaggedRight

        \color{jFrametitleFGColor}

        The content of this lecture is made available under a Creative
        Commons licence (\textsc{cc}), giving the author the credits
        (\textsc{by}) and putting an obligation to share on the same terms
        (\textsc{sa}). Figures and diagrams included in the lecture are
        authored by Paweł Węgrzyn et~al., and are available under the same
        license unless indicated otherwise.\\ The presentation uses the
        Beamer Jagiellonian theme based on Matthias Vogelgesang’s
        Metropolis theme, available under license \LaTeX{} Project
        Public License~1.3c at: \colorhref{https://github.com/matze/mtheme}
        {https://github.com/matze/mtheme}.

        Typographic design: Iwona Grabska-Gradzińska;
        \LaTeX{} Typesetting: Kamil Ziemian \\
        Proofreading: Wojciech Palacz,
        Monika Stawicka \\
        3D Models: Dariusz Frymus, Kamil Nowakowski \\
        Figures and charts: Kamil Ziemian, Paweł Węgrzyn,
        Wojciech Palacz \\
        Film editing: Agencja Filmowa Film \& Television Production~--
        Zbigniew Masklak

      \end{flushleft}

    \end{textblock}

  \end{frame}





  \begin{frame}[standout]


    \begingroup

    \color{jFrametitleFGColor}

    #1

    \endgroup

  \end{frame}
}



\newcommand{\GeometryThreeDTwoSpecialEndingSlidesVideoVerTwoEN}[1]{%
  \begin{frame}[standout]

    \begin{textblock}{11}(1,0.7)

      \begin{flushleft}

        \mdseries

        \footnotesize

        \color{jFrametitleFGColor}

        This content was created as part of a project co-financed by the
        European Union within the framework of the European Social Fund
        POWR.03.05.00-00-Z309/17-00.

      \end{flushleft}

    \end{textblock}





    \begin{textblock}{10}(0,2.2)

      \tikz \fill[color=jBackgroundStyleLight] (0,0) rectangle (12.8,-1.5);

    \end{textblock}


    \begin{textblock}{3.2}(0.7,2.45)

      \includegraphics[scale=0.3]{\FundingLogoColorPictureEN}

    \end{textblock}


    \begin{textblock}{2.5}(4.15,2.5)

      \includegraphics[scale=0.2]{\JULogoColorPictureEN}

    \end{textblock}


    \begin{textblock}{2.5}(6.35,2.4)

      \includegraphics[scale=0.1]{\ZintegrUJLogoColorPictureEN}

    \end{textblock}


    \begin{textblock}{4.2}(8.4,2.6)

      \includegraphics[scale=0.3]{\EUSocialFundLogoColorPictureEN}

    \end{textblock}





    \begin{textblock}{11}(1,4)

      \begin{flushleft}

        \mdseries

        \footnotesize

        \RaggedRight

        \color{jFrametitleFGColor}

        The content of this lecture is made available under a Creative
        Commons licence (\textsc{cc}), giving the author the credits
        (\textsc{by}) and putting an obligation to share on the same terms
        (\textsc{sa}). Figures and diagrams included in the lecture are
        authored by Paweł Węgrzyn et~al., and are available under the same
        license unless indicated otherwise.\\ The presentation uses the
        Beamer Jagiellonian theme based on Matthias Vogelgesang’s
        Metropolis theme, available under license \LaTeX{} Project
        Public License~1.3c at: \colorhref{https://github.com/matze/mtheme}
        {https://github.com/matze/mtheme}.

        Typographic design: Iwona Grabska-Gradzińska;
        \LaTeX{} Typesetting: Kamil Ziemian \\
        Proofreading: Wojciech Palacz,
        Monika Stawicka \\
        3D Models: Dariusz Frymus, Kamil Nowakowski \\
        Figures and charts: Kamil Ziemian, Paweł Węgrzyn,
        Wojciech Palacz \\
        Film editing: IMAVI -- Joanna Kozakiewicz, Krzysztof Magda, Nikodem
        Frodyma

      \end{flushleft}

    \end{textblock}

  \end{frame}





  \begin{frame}[standout]


    \begingroup

    \color{jFrametitleFGColor}

    #1

    \endgroup

  \end{frame}
}



\newcommand{\GeometryThreeDTwoSpecialEndingSlidesVideoVerThreeEN}[1]{%
  \begin{frame}[standout]

    \begin{textblock}{11}(1,0.7)

      \begin{flushleft}

        \mdseries

        \footnotesize

        \color{jFrametitleFGColor}

        This content was created as part of a project co-financed by the
        European Union within the framework of the European Social Fund
        POWR.03.05.00-00-Z309/17-00.

      \end{flushleft}

    \end{textblock}





    \begin{textblock}{10}(0,2.2)

      \tikz \fill[color=jBackgroundStyleLight] (0,0) rectangle (12.8,-1.5);

    \end{textblock}


    \begin{textblock}{3.2}(0.7,2.45)

      \includegraphics[scale=0.3]{\FundingLogoColorPictureEN}

    \end{textblock}


    \begin{textblock}{2.5}(4.15,2.5)

      \includegraphics[scale=0.2]{\JULogoColorPictureEN}

    \end{textblock}


    \begin{textblock}{2.5}(6.35,2.4)

      \includegraphics[scale=0.1]{\ZintegrUJLogoColorPictureEN}

    \end{textblock}


    \begin{textblock}{4.2}(8.4,2.6)

      \includegraphics[scale=0.3]{\EUSocialFundLogoColorPictureEN}

    \end{textblock}





    \begin{textblock}{11}(1,4)

      \begin{flushleft}

        \mdseries

        \footnotesize

        \RaggedRight

        \color{jFrametitleFGColor}

        The content of this lecture is made available under a Creative
        Commons licence (\textsc{cc}), giving the author the credits
        (\textsc{by}) and putting an obligation to share on the same terms
        (\textsc{sa}). Figures and diagrams included in the lecture are
        authored by Paweł Węgrzyn et~al., and are available under the same
        license unless indicated otherwise. \\ The presentation uses the
        Beamer Jagiellonian theme based on Matthias Vogelgesang’s
        Metropolis theme, available under license \LaTeX{} Project
        Public License~1.3c at: \colorhref{https://github.com/matze/mtheme}
        {https://github.com/matze/mtheme}.

        Typographic design: Iwona Grabska-Gradzińska;
        \LaTeX{} Typesetting: Kamil Ziemian \\
        Proofreading: Leszek Hadasz, Wojciech Palacz,
        Monika Stawicka \\
        3D Models: Dariusz Frymus, Kamil Nowakowski \\
        Figures and charts: Kamil Ziemian, Paweł Węgrzyn,
        Wojciech Palacz \\
        Film editing: Agencja Filmowa Film \& Television Production~--
        Zbigniew Masklak \\
        Film editing: IMAVI -- Joanna Kozakiewicz, Krzysztof Magda, Nikodem
        Frodyma


      \end{flushleft}

    \end{textblock}

  \end{frame}





  \begin{frame}[standout]


    \begingroup

    \color{jFrametitleFGColor}

    #1

    \endgroup

  \end{frame}
}











% ------------------------------------------------------
% BibLaTeX
% ------------------------------------------------------
% Package biblatex, with biber as its backend, allow us to handle
% bibliography entries that use Unicode symbols outside ASCII.
\usepackage[
language=polish,
backend=biber,
style=alphabetic,
url=false,
eprint=true,
]{biblatex}

\addbibresource{Systemy-operacyjne-Bibliography.bib}





% ------------------------------------------------------
% Importing packages, libraries and setting their configuration
% ------------------------------------------------------





% ------------------------------------------------------
% Local packages
% ------------------------------------------------------
% Local configuration of this particular presentation
\usepackage{./Local-packages/local-settings}

\usepackage{./Local-packages/PGF-TikZ-Arrows-styles}

\usepackage{./Local-packages/PGF-TikZ-Diagram-styles}










% ####################################################################
\title{Systemy operacyjne}
\subtitle{3.~Pewne informacje o~systemach operacyjnych}

\author{Kamil Ziemian \\
  \email}


% \date{}
% ####################################################################










% ####################################################################
% Beginning of the document
\begin{document}
% ####################################################################





% ################################################
% Number of chars: 43k+, 8k+,
% Text is adjusted to the left and words are broken at the end of the line.
\RaggedRight
% ################################################





% ################################################
\maketitle
% ################################################





% ################################################
\begin{frame}
  \frametitle{Spis treści}


  \tableofcontents

\end{frame}
% ################################################










% ################################################
\section{Systemy operacyjne wokół nas}
% ################################################


% ######################################
\begin{frame}
  \frametitle{Po co istnieją systemy operacyjne i~czemu jest ich
    tak wiele?}


  Jeśli rozejrzymy~się po świecie, dostrzeżemy jak wielka liczba różnych
  systemów operacyjnych (os, ang. \textit{Operating System}) jest dziś
  używana, choć normalny użytkownik
  styka~się bezpośrednie tylko z~kilkoma z~nich. Dlaczego jest ich tak
  wiele? Spróbujmy wskazać kilka przyczyn tego stanu rzeczy.

  Pierwszą kwestią o~której należy wspomnieć jest to, że~system operacyjny
  służy do zarządzania sprzętem komputerowym, różne rodzaje komputerów
  wymagają różnych systemów operacyjnych. Inny system operacyjny potrzebny
  jest na komputerze stacjonarnym, inny na smartfonie. To jedna przyczyna
  tej różnorodności.

  Drugą kwestią jest to, że~różni ludzie mają różne pomysły na temat tego
  jak system operacyjny powinien działać. Więc nawet na ten sam typ
  komputera ludzie napiszą różne systemy operacyjne, które według nich
  lepiej wykonujące swoją pracę, niż inne.

\end{frame}
% ######################################





% ######################################
\begin{frame}
  \frametitle{Po co istnieją systemy operacyjne i~czemu jest ich
    tak wiele?}


  Po trzecie, nawet jeśli rdzeń systemu operacyjnego jest gotowy, to
  istnieje milion rzeczy dokoła niego, które można zmienić. By użyć bardzo
  prostych przykładów, jak ma wyglądać jego środowisko graficzne,
  jakie programy do edycji zdjęć mają być na nim dostępne, jaki program
  ma odtwarzać na nim muzykę, etc.? Taka właśnie sytuacja ma miejsce
  w~systemie GNU/Linux, który posiada wiele dystrybucji. Do zagadnienie
  tego krótko powrócimy w~przyszłości.

  Przejdziemy teraz do krótkiego przegląd systemów operacyjnych używanych
  w~XXI wieku. To oczywiście bardzo obszerny temat, któremu można by
  poświęcić wiele godzin. Poniższy przegląd ma na celu tylko wstępne
  uporządkowanie pewnych rzecz, które mogli już Państwo usłyszeć wcześniej.

\end{frame}
% ######################################





% ######################################
\begin{frame}
  \frametitle{Jakie systemy operacyjne są dziś używane?}


  Obecnie na komputerach osobistych dominują trzy rodziny systemów
  operacyjnych. Są to:
  \begin{itemize}

  \item \colorhref{https://www.gnu.org/gnu/linux-and-gnu.en.html}{GNU/Linux}
    i~\colorhref{https://en.wikipedia.org/wiki/Linux\_kernel}{Linux};

  \item \colorhref{https://en.wikipedia.org/wiki/MacOS}{macOS};

  \item \colorhref{https://en.wikipedia.org/wiki/Microsoft\_Windows}
    {Microsoft Windows}.

  \end{itemize}

  Gdy chodzi o~smartfony, to istnieją co~najmniej dwie duże rodziny systemów
  operacyjnych dla nich:
  \begin{itemize}

  \item
    \colorhref{https://en.wikipedia.org/wiki/Android\_(operating\_system)}
    {Android};

  \item \colorhref{https://en.wikipedia.org/wiki/IOS}{iOS}.

  \end{itemize}
  Niemniej nie jesteśmy znawcami smartfonów, więc nie będziemy zbytnio
  poruszali tej tematyki.


\end{frame}
% ######################################





% ######################################
\begin{frame}
  \frametitle{Jakie systemy operacyjne są dziś używane?}


  Jeśli chodzi o~inne systemy operacyjne dla komputerów osobistych to
  można wymienić je~bez końca:
  \colorhref{https://en.wikipedia.org/wiki/Fire\_OS}{Fire~OS},
  \colorhref{https://en.wikipedia.org/wiki/FreeBSD}{FreeBSD},
  FreeDOS, Haiku,
  HarmonyOS, HeliOS, Inferno, MINIX, OpenHarmony, OpenSolaris,
  Phantom~OS, Plan~9 from Bell Labs, ReactOS, Redox, Thesueus~OS
  (interesujący projekt), Visopsys, etc.

  Należy pamiętać, że~komputery osobiste, to tylko jeden z~typów
  komputerów, których obecnie używamy. Z~innych typów komputerów można
  wymienić smartfony, komputery mainframowe, serwery, komputery wbudowane
  w~mikrofalówki,~etc.

  Najwięcej różnych systemów operacyjnych jest zapewne na komputery
  osobiste, bo są one powszechnie dostępne, a~tworzenie dla nich
  oprogramowania jest relatywnie proste. Stąd jeśli kilku dobrych
  programistów uważa, że~warto stworzyć dla nich nowy system operacyjny,
  to niedługo będziemy mogli dopisać jeszcze jedną pozycję do~podanej
  wcześniej listy.

\end{frame}
% ######################################





% ######################################
\begin{frame}
  \frametitle{Komputery i~ich systemy operacyjne}


  Należy mieć świadomość, że~jeden system operacyjny może pracować na kilku
  typach komputerów, czego doskonałym przykładem jest Linux. Z~drugiej
  strony, różne typy komputery mają różne możliwości i~potrzeby, inaczej
  działa komputer mainframe zajmujący cały pokój, inaczej smartfon
  mieszczący~się w~kieszeni. Dlatego też często na komputery różnych typów,
  pisze~się systemy operacyjne przeznaczone specjalnie dla nich, dostosowane
  do ich~specyfiki.

  Pozwolimy sobie rzucić krótko okiem na różne typy komputerów i~dostępne
  na nich systemy operacyjne.

  \begin{itemize}

  \item Systemy operacyjne komputerów mainframowych. Ten rodzaj komputerów
    spotyka~się zwykle w~centrach obliczeniowych różnorakiego
    przeznaczenia. Obsługiwane są przez specjalne dystrybucje Linuxa,
    jak~również systemy operacyjne stworzone specjalnie dla nich,
    jak~\colorhref{https://en.wikipedia.org/wiki/Z/OS}{z/OS}.

  \end{itemize}

\end{frame}
% ######################################





% ######################################
\begin{frame}
  \frametitle{Komputery i~ich systemy operacyjne}


  \begin{itemize}

  \item Systemy operacyjne serwerów. Tutaj również Linux jest popularny.

  \item Systemy operacyjne komputerów osobistych, nim właśnie jest
    poświęcony ten przedmiot.

  \item Systemy operacyjne smartfonów, jak Android czy iOS. Należy
    dodać, że~Android sam jest odmianą Linuxa.

  \item Wbudowane systemy operacyjne. Są to systemy operacyjne obsługujące
    pralki, kuchenki mikrofalowe, samochody, etc.

  \item Systemy operacyjne kart elektronicznych. Temat na zupełnie
    osobne zajęcia.

  \item Systemy operacyjne węzłów sensorowych. Przykładem węzła sensorowego
    są czujniki przeciwpożarowe, układy do~mierzenia temperatury, ilości
    opadów, etc.

  \end{itemize}

\end{frame}
% ######################################





% ######################################
\begin{frame}
  \frametitle{Komputery i~ich systemy operacyjne}


  Ponieważ życie nie może być proste, podział systemów operacyjnych
  ze~względu na to na jakim sprzęcie działają, nie jest jedynym
  istniejącym. Systemy operacyjne można też podzielić na monolityczne,
  wielowarstwowe, czasu rzeczywistego, etc.

  \alert{Nie} muszą się Państwo uczyć powyższej przedstawionych listy
  systemów operacyjnych. Podaliśmy je, by mieli Państwo świadomość,
  że~systemy operacyjne to niełatwa część informatyki i~należy odczuwać
  zdrową porcję szacunku wobec ludzi którzy te systemy tworzą. Pamiętając,
  że~ci sami ludzie całkiem dużo rzeczy mocno spaprali.

  Zwróćmy uwagę na jeszcze jedną rzecz. Tak jak nasz świat jest pełen
  komputerów, tak jest pełen systemów operacyjnych. Bez wielkiej przesady
  można powiedzieć, że~bez systemów operacyjnych setki milionów komputerów
  stałby~się dla większości ludzi bezużyteczną stertą złomu.
  \colorhref{https://en.wikipedia.org/wiki/Windows\_11}{Windows~$11$?}

\end{frame}
% ######################################





% ######################################
\begin{frame}
  \frametitle{Maszyny Rube Goldbergera}


  Informatyka w~ogólności, a~systemy operacyjne w~szczególności, pełne są
  maszyn
  \colorhref{https://en.wikipedia.org/wiki/Rube\_Goldberg\_machine}{Rube
    Goldbergera}. Państwo już na pewno nie raz
  widzieli taką maszynę, a~jedną z~nich można obejrzeć
  \colorhref{https://www.youtube.com/watch?v=vn-g1Mn2\_3g}{tutaj}.
  Dlatego nie zawsze warto~się zastanawiać, czemu coś w~systemie operacyjnym
  działa w~dziwaczny sposób. Działa w~taki sposób, bo kiedyś ktoś nie miał
  lepszego pomysłu i~tak już zostało.

\end{frame}
% ######################################










% ################################################
\section{Dlaczego na tych zajęciach używamy systemu
  GNU/Linux?}
% ################################################


% ######################################
\begin{frame}
  \frametitle{Dlaczego na tych zajęciach używamy systemu
    GNU/Linux?}


  Przypomnijmy jaka jest różnica między GNU/Linuxem, a~Linuxem.
  Najważniejsza część systemu operacyjnego nosi nazwę jądra
  (ang.~\textit{kernel}), kontroluje ono działanie sprzętu komputerowego,
  takiego jak procesor, dysk, ekran, drukarka,~etc. Każdy system
  operacyjny posiada właściwe jemu jądro. Jednak nikt
  normalny nie używa systemu operacyjnego składającego~się tylko z~jądra.
  Potrzebujemy też takich rzeczy jak program do pracy z~pulpitem, o~których
  istnieniu zwykle nawet nie myślimy.

  Przez „system operacyjny Linux” rozumiem każdy system operacyjny,
  który używa jądra Linuxa. Przez „GNU/Linuxa” rozumiem system operacyjny,
  który zawiera jądro Linuxa i~do~wykonywania tysięcy rzeczy, których jądro
  nie obsługuje bezpośrednio, używa zestawu programów dostarczonych przez
  projekt GNU. Przypomnijmy, że~„GNU” to rekurencyjny akronim „GNU is Not
  a~Unix”.

\end{frame}
% ######################################





% ######################################
\begin{frame}
  \frametitle{Dlaczego na tych zajęciach używamy systemu
    GNU/Linux?}


  Dlaczego na tych zajęciach zdecydowaliśmy~się korzystać z~systemu
  GNU/Linux? W~zasadzie każdy typ systemów operacyjnych (systemy
  smartfonów, komputerów osobistych, komputerów mainfram, etc.) to osobna
  gałąź informatyki, więc niezależnie od tego, reprezentanta którego
  z~tych typów byśmy nie wybrali, dokonamy pewnego zawężenia tematyki
  naszego przedmiotu. Trzeba jednak podkreślić, że~jeśli dobrze poznamy
  systemy operacyjne dla komputerów osobistych, to nauka systemów
  operacyjny smartfonów będzie łatwiejsza. Łatwiejsza, ale~pewnych rzeczy
  trzeba~się będzie uczyć od nowa. Co więc przemawia za~wybraniem
  GNU/Linuxa?

  Zacznijmy od czysto praktycznych względów. System ten jest darmowy, więc
  jeśli ktoś posiada komputer, może go zainstalować na nim, nie płacą za
  niego nic. Zainstalowanie go wymagać będzie prawie na pewno pewnej
  dozy wiedzy i~umiejętności, ale dla nikogo nie stanowi bariery finansowej.

\end{frame}
% ######################################





% ######################################
\begin{frame}
  \frametitle{Dlaczego na tych zajęciach używamy systemu
    GNU/Linux?}


  Linux jest systemem posiadający wielkie praktyczne znaczenie, działa
  na setkach milionów urządzeń, od smartfonów po superkomputery, więc
  umiejętność jego używania może zostać wykorzystana do rozwiązywania
  konkretnych życiowych problemów. Linux jest więc antytezą akademickiego
  systemu operacyjnego, który służy tylko do tego, by prowadzić w~oparciu
  o~niego zajęcia na jakiejś uczelni.

  W~końcu, GNU/Linux jest projektem o~otwartym kodzie źródłowym, więc jeśli
  tylko chcemy i~jesteśmy odpowiednio twardzi, możemy zajrzeć do~każdego
  jego zakamarka i~dowiedzieć~się o~nim \alert{wszystkiego}. Nigdzie
  w~nauce nie napotkamy bariery, finansowej czy prawnej. Barierą
  może być tylko trudność zrozumienia danego zagadnienia.

\end{frame}
% ######################################










% ################################################
\section{Ile tych GNU/Linuksów jest?}
% ################################################


% ######################################
\begin{frame}
  \frametitle{antiX, Astra Linux, Canaima, ile tego jeszcze?}


  W~przeciwieństwie do choćby systemu Windows~$10$, który jest jednym
  systemem operacyjnym, GNU/Linux jest raczej pstrokatym zbiorowiskiem
  przeróżnych swoich odmian. Żeby to wyjaśnić, zacznijmy od~komentarza
  odnośnie Windowsa~$10$. System ten występuje w~kilku wersjach, w~tym
  sensie, że~firma Microsoft wydaje kolejne updejty ;) tego systemu, jednak
  to~wciąż jest ten sam system operacyjny, w~którym zmieniono kilka trzeci-
  i~czwartorzędnych rzeczy.

  Aby było jasne, updaty często zawierają bardzo ważne z~punktu widzenia
  użytkownika rzeczy, jak na przykład kod ograniczający możliwość
  zhakowania urządzeni, ale z~punktu widzenia działania systemu
  operacyjnego, są to kwestie trzeciorzędnej natury.

  W~przeciwieństwie do tego, GNU/Linux występuje w~rozmaitych odmianach,
  zwanych \alert{dystrybucjami}, których liczba jest trudna do ogarnięcia.
  Wymieńmy niektóre z~nich.
  \colorhref{https://antixlinux.com/}{antiX},
  \colorhref{https://astralinux.ru/}{Astra Linux},
  \colorhref{https://archlinux.org/}{Arch Linux},
  \colorhref{https://canaima.softwarelibre.gob.ve/}{Canaima},
  \colorhref{https://www.debian.org/}{Debian},
  \colorhref{https://www.gentoo.org/}{Gentoo},
  \colorhref{https://ubuntu.com/}{Ubuntu},\ldots, lista ta zdaje~się
  nie mieć końca.

\end{frame}
% ######################################





% ######################################
\begin{frame}
  \frametitle{antiX, Astra Linux, Canaima, ile tego jeszcze?}


  Dla osób przytłoczonych liczbą tych wersji systemu GNU/Linux, warto
  dodać, że~choć liczba samych dystrybucji jest ogromna, to można ją
  podzielić na dosłownie kilka rodzin, jak choćby systemy wywodzące~się od
  dystrybucji
  \colorhref{https://en.wikipedia.org/wiki/Debian}{Debian}. Ponieważ
  dystrybucje wewnątrz jednej rodziny~są do siebie podobne, czyni to całą
  sytuację możliwą do ogarnięcia.

  Wszystkie te dystrybucje używają oczywiście jądra Linuxa, więc
  trafne jest stwierdzenie, że~wszystkie są odmianą jednego systemu
  operacyjnego. Dystrybucje GNU/Linuxa, z~definicji, używają też tych
  samych narzędzi projektu GNU.

  Przytłaczająca większość dystrybucji Linuxa jest darmowa. Wchodzisz
  na odpowiednią stronę internetową i~pobierasz wszystko czego chcesz.
  Istnieje kilka płatnych dystrybucji Linuxa, z~których najsłynniejszą jest
  zapewne
  \colorhref{https://en.wikipedia.org/wiki/Red\_Hat\_Enterprise\_Linux}
  {Red Hat Enterprise Linux}, w~skrócie RHEL, tworzony przez firmę
  \colorhref{https://en.wikipedia.org/wiki/Red\_Hat}{Red Hat Inc.}.

\end{frame}
% ######################################




% ######################################
\begin{frame}
  \frametitle{antiX, Astra Linux, Canaima, ile tego jeszcze?}


  Dla osób przytłoczonych liczbą tych wersji systemu GNU/Linux, warto
  dodać, że~choć liczba samych dystrybucji jest ogromna, to można ją
  podzielić na dosłownie kilka rodzin, jak choćby systemy wywodzące~się od
  dystrybucji
  \colorhref{https://en.wikipedia.org/wiki/Debian}{Debian}. Ponieważ
  dystrybucje wewnątrz jednej rodziny~są do siebie podobne, czyni to całą
  sytuację możliwą do ogarnięcia.

  Wszystkie te dystrybucje używają oczywiście jądra Linuxa, więc
  trafne jest stwierdzenie, że~wszystkie są odmianą jednego systemu
  operacyjnego. Dystrybucje GNU/Linuxa, z~definicji, używają też tych
  samych narzędzi projektu GNU.

  Przytłaczająca większość dystrybucji Linuxa jest darmowa. Wchodzisz
  na odpowiednią stronę internetową i~pobierasz wszystko czego chcesz.
  Istnieje kilka płatnych dystrybucji Linuxa, z~których najsłynniejszą jest
  zapewne
  \colorhref{https://en.wikipedia.org/wiki/Red\_Hat\_Enterprise\_Linux}
  {Red Hat Enterprise Linux}, w~skrócie RHEL, tworzony przez firmę
  \colorhref{https://en.wikipedia.org/wiki/Red\_Hat}{Red Hat Inc.}.

\end{frame}
% ######################################





% ######################################
\begin{frame}
  \frametitle{Fundacja Linuxa}


  W~tym kontekście warto wspomnieć pewne osobliwości związane z~Linuxem
  i~światem w~którym żyjemy. Rozwojem jądra tego systemu zajmuje~się
  \colorhref{https://en.wikipedia.org/wiki/Linux\_Foundation}
  {Fundacja Linuxa}, która jest zarejestrowana jako organizacja non-profit.
  Mimo tego, można znaleźć informacje, że~w~roku $2024$ przychody tej
  fundacji wyniosły $300$~milionów dolarów, z~czego mniej niż $3\%$~(!)
  przeznaczono na rozwój jądra Linuxa. Jak słusznie zwrócił uwagę kanał
  The~Lunduke Journal, tego typu wyniki finansowe mogłyby zawstydzić
  niejedną firmę, której celem jest generowanie przychodów
  \parencite{TLJ-Linux-Foundation-Drops-Linux-Spending-to-Historic-ETC-Ver-2024}.

  Skąd ta Fundacja bierze taką ilość pieniędzy? Według The~Lunduke Journal,
  $125$ milionów dolarów pochodziło od korporacji wspierających fundację.
  Co jest ważne w~tym kontekście, najwięksi darczyńcy otrzymują w~ramach
  tego, co po angielsku nosi nazwę \textit{corporate membership}, prawo
  do~zasiadania w~zarządzie fundacji (ang.~\textit{board~of directors}).

\end{frame}
% ######################################





% ######################################
\begin{frame}
  \frametitle{Fundacja Linuxa}


  W~roku $2026$,
  w~\colorhref{https://www.linuxfoundation.org/about/leadership}
  {zarządzie fundacji Linuxa} zasiadali przedstawiciele firm Sony, Red
  Hat, Huawei, Renesas, Fujitsu, Qualcomm Technologies,~Inc., Meta,
  Hitachi, Panasonic, Microsoft, Samsung, Ericsson, NEC, Intel i~Oracle.

\end{frame}
% ######################################





% ##################
\begin{frame}
  \frametitle{Standaryzacja}


  Obecnie, piszemy to w~$2025$, GNU/Linux jest dobrze zestandaryzowany
  na~poziomie swoich podstawowych zasad działania i~funkcji, które oferuje
  użytkownikom. W~istocie, wiele własności GNU/Linuxa sięga rozwijanego od
  około $1970$~roku systemu operacyjnego
  \colorhref{https://en.wikipedia.org/wiki/Unix}{\textsc{unix}} i~zestawu
  standardów dla tego systemu znanej jako
  \colorhref{https://en.wikipedia.org/wiki/POSIX}{\textsc{posix}}
  (ang.~\textit{Portable Operating System Interface for Unix}). Dla
  kontekstu, standardy \textsc{posix} zaczęły być wydawane przez Instytut
  Inżynierów Elektryków i~Elektroników (\textsc{ieee},
  ang.~\textit{Institute of Electrical and Electronics Engineers})
  w~roku~$1988$.

  Dlatego też większość rzeczy, które poznają Państwo na~tym kursie powinny
  działać w~$99\%$ dystrybucji systemu GNU/Linux i~w~$100\%$ dystrybucji
  które są używane w~poważnych firmach i~projektach. Będziemy~się też
  starać wyraźnie zaznaczać, czym mogą~się różnić poszczególne dystrybucje
  GNU/Linuxa.

\end{frame}
% ##################










% ################################################
\section{Zmienne w~BASHu}
% ################################################



% ######################################
\begin{frame}
  \frametitle{Zmienne w~BASHu}


  Koncepcje zmiennej na pewno znają Państwo już z~innych języków
  programowania i~tłumaczenie ich użyteczności nie powinno być konieczne.
  Jednak zmienne w~BASHu, jak chyba wszystko w~BASHu, mają kilka osobliwych
  cech do których musimy przywyknąć.

  Zacznijmy od zdefiniowania zmiennej w~BASHu. \\
  \texttt{\$ var1="Jestem zmienną."} \\
  Powyższa linia definiuje zmienną o~nazwie \texttt{var1}, która zawiera
  string „\texttt{Jestem zmienną.}”. Proszę bardzo uważać
  na~\alert{białe znaki} w~przedstawionym poleceniu. BASH nie zna słowa
  litość, a~gdy chodzi o~definiowanie zmiennych jest jeszcze bardziej
  bezlitosny, niż normalnie. Jeśli to w~ogóle możliwe.

  Przykładowo, linia \\
  \texttt{\$ var1 = "Jestem zmienną."} \\
  jest \alert{niepoprawna}! Bo czemu życie miałoby być proste?

\end{frame}
% ######################################





% ######################################
\begin{frame}
  \frametitle{BASH nie zna słowa litość. Odcinek kolejny}


  Jeśli będziemy próbowali uruchomić polecenie \\
  \texttt{\$ var1 = "Jestem zmienną."} \\
  to zobaczymy komunikat o~błędzie, który może wyglądać w~następujący
  sposób. \\
  \texttt{> Nie znaleziono polecania 'var1', czy chodziło~o: \ldots} \\
  Jak zwykle, konkretna postać komunikatu o~błędzie zależy od używanej
  dystrybucji GNU/Linuxa.

  Problem bierze~się stąd, że~normalnie BASH oczekuje, iż~pierwsze słowo
  które napotka w~danej linii jest nazwą jakiegoś polecenia. W~przypadku
  polecania podanego wyżej, pierwszym tym słowem jest „\texttt{var1}” i~BASH
  narzeka, że~polecenia o~tej nazwie nie zna. Bo czemu życia miałoby być
  proste?

\end{frame}
% ######################################





% ######################################
\begin{frame}
  \frametitle{BASH nie zna słowa litość. Odcinek kolejny}


  Dopiero zapisanie polecenia \\
  \texttt{\$ var1="Jestem zmienną."} \\
  bez spacji, zmusi BASHa do przeczytania całej jego treści
  i~zorientowania~się, że~jest to instrukcja przypisania wartości
  do~zmiennej.

  Operację przypisania wartości do zmiennej, też nazywamy poleceniem,
  bo z~punktu widzenia BASHa, nim właśnie jest. Przypomnijmy, że~w~naszym
  kontekście, „polecenie” to dowolny kod, który poprawnie opisuje działanie
  jakie może wykonać BASH (lub inna powłoka).

  Byłoby pięknie, gdyby na tym problemy ze zmiennymi~się kończyły, ale
  przecież życie nie może być proste. Powiedzmy, że~chcemy odczytać
  zawartość naszej zmiennej \texttt{var1}. Wydawałoby~się, że~powinien
  zadziałać kod \\
  \texttt{\$ echo "var1"}

\end{frame}
% ######################################





% ######################################
\begin{frame}
  \frametitle{BASH nie zna słowa litość. Odcinek kolejny}


  Oto co~się naprawdę stanie: \\
  \texttt{\$ echo "var1"} \\
  \texttt{> var1} \\
  Standardowo w~każdym języku programowania, każde wystąpienie zmiennej
  zastępowane jest jej wartością. Jednak nie w~BASHu, bo czemu życie miałoby
  być proste?

  Aby zobaczyć wartość przechowywaną przez zmienną w~BASHu, należy jej
  nazwę poprzedzić znakiem „\$”: \\
  \texttt{\$ echo "\$var1"} \\
  \texttt{> Jestem zmienną.} \\
  Ze znanych nam języków programowania jedynie BASH i~inne języki podobne do
  niego języki powłok mają taką dziwaczną własność.

\end{frame}
% ######################################





% ######################################
\begin{frame}
  \frametitle{BASH nie zna słowa litość. Odcinek kolejny}


  By to lepiej zilustrować, przetestujmy poniższe polecenie: \\
  \texttt{\$ echo "var1: |\$var1|."} \\
  \texttt{> var1: |Jestem zmienną.|.} \\
  W~przykładzie tym użyliśmy kresek pionowych, by wyraźniej zaznaczyć,
  co~jest przechowywane przez zmienną \texttt{var1}.

  Na tym oczywiście nasze problemy, nie mogę się kończyć, o~nie. Przyjmijmy,
  że~zmienna \texttt{var2} jest obecnie niezdefiniowana. W~takiej sytuacji,
  polecenie \\
  \texttt{\$ echo "var2: |\$var2|."} \\
  powinno zgłosić błąd.

  Powinno, ale to by znaczyło, że~BASH zna pojęcie litości. Oto co~się
  naprawdę stanie: \\
  \texttt{\$ echo "var2: |\$var2|."} \\
  \texttt{> var2: ||.} \\

\end{frame}
% ######################################





% ######################################
\begin{frame}
  \frametitle{BASH nie zna słowa litość. Odcinek kolejny}


  \texttt{\$ echo "var2: |\$var2|."} \\
  \texttt{> var2: ||.} \\
  Z~jakiegoś powodu, gdy BASH napotka nieznaną sobie zmienną, to
  \alert{zamiast} zgłosić błąd, jak powinien, zwraca jako jej wartość
  \alert{pusty string}! Powinno być jasne, że~jest to źródło trudnej
  do~ogarnięcie liczby błędów. Bo czemu życie miałoby być proste?

  Morał z~tej historii jest taki. Jeśli w~pracy z~BASHem używasz
  zmiennych i~coś nie działa on jak trzeba, to w~pierwszym rzędzie sprawdź,
  czy nie zrobiłeś gdzieś literówki w~nazwie zmiennej.

  Jeśli wpisując nazwę jakiejś zmiennej zrobiłeś gdzieś literówkę,
  to~najprawdopodobniej wywołałeś nieistniejącą zmienną, więc zamiast tego
  co powinno tam być, w~tym miejscu znajduje~się teraz pusty string.
  I~to niszczy wszystko.

\end{frame}
% ######################################





% ######################################
\begin{frame}
  \frametitle{BASH nie zna słowa litość. Odcinek kolejny}


  \texttt{\$ echo "var2: |\$var2|."} \\
  \texttt{> var2: ||.} \\
  Z~jakiegoś powodu, gdy BASH napotka nieznaną sobie zmienną, to
  \alert{zamiast} zgłosić błąd, jak powinien, zwraca jako jej wartość
  \alert{pusty string}! Powinno być jasne, że~jest to źródło trudnej
  do~ogarnięcie liczby błędów. Bo czemu życie miałoby być proste?

  Morał z~tej historii jest taki. Jeśli w~pracy z~BASHem używasz
  zmiennych i~coś nie działa on jak trzeba, to w~pierwszym rzędzie sprawdź,
  czy nie zrobiłeś gdzieś literówki w~nazwie zmiennej.

  Jeśli wpisując nazwę jakiejś zmiennej zrobiłeś gdzieś literówkę,
  to~najprawdopodobniej wywołałeś nieistniejącą zmienną, więc zamiast tego
  co powinno tam być, w~tym miejscu znajduje~się teraz pusty string.
  I~to niszczy wszystko.

\end{frame}
% ######################################





% ######################################
\begin{frame}
  \frametitle{Korzystanie ze~zmiennych w~BASHu}


  Jak już widzieliśmy, używanie zmienne w~BASHu nie należy do najprostszych
  rzeczy na~świecie. Jednak zmienne są tak bardzo użyteczne, że~nie należy
  ich unikać. Jak mawiają w~świecie gier wideo, wszystko co musisz zrobić
  to~„Git gud!”.

  BASH cały czas korzysta z~wielu zmiennych, które noszą nazwę „zmiennych
  środowiskowych” (ang. \textit{environment variables}). Zmienne te służą
  do konfigurowania działania BASHa. Przykładowo, skąd BASH wie, gdzie
  jest nasz katalog domowy? Dlaczego polecenie \texttt{cd $\sim$}
  przenosi nas właśnie do tego katalogu, a~nie jakiegoś innego?

  Za każdym razem, gdy BASH potrzebuje skorzystać, ze ścieżki do katalogu
  domowego, jak w~poleceniu \texttt{cd $\sim$}, sprawdza on zawartość
  zmiennej \texttt{HOME}. Możemy to sprawdzić, pisząc \\
  \texttt{\$ echo "\$HOME"} \\
  \texttt{> /home/mojeImie}

\end{frame}
% ######################################





% ######################################
\begin{frame}
  \frametitle{Korzystanie ze~zmiennych w~BASHu}


  \texttt{\$ echo "\$HOME"} \\
  \texttt{> /home/mojeImie} \\
  Adres czyjegoś katalogu domowego zależy od tego, jak konkretnie
  ustawione~są katalogi na danym komputerze, na co nie ma żadnej reguły.
  Nie należy więc~się spodziewać, że~na Państwa komputerze ścieżka do
  katalogu domowego będzie wyglądać jak podana wyżej. Może nawet nie być
  zbyt do niej podobna.

\end{frame}
% ######################################





% ################################################
\section{Zdalne logowanie~się do GNU/Linuxa z~Windowsa}
% ################################################







% ################################################
\appendix
% ################################################





% ######################################
\EndingSlide{Dziękujmy! Pytania?}
% ######################################










% ####################################################################
% ####################################################################
% Bibliography

\printbibliography





% ####################################################################
% End of the document

\end{document}
